\documentclass[10pt,a4paper]{amsart}
\usepackage{amsmath,amssymb,amsthm,graphicx}
\usepackage{verbatim}
\newenvironment{reference}{\comment}{\endcomment}
%\newenvironment{reference}{}{}
\newenvironment{slogan}{\comment}{\endcomment}
\newenvironment{history}{\comment}{\endcomment}

% For commutative diagrams we use Xy-pic
\usepackage[all]{xy}

% We use 2cell for 2-commutative diagrams.
\xyoption{2cell}
% Theorem environments.
%
\theoremstyle{plain}
\newtheorem{theorem}[subsection]{Theorem}
\newtheorem{proposition}[subsection]{Proposition}
\newtheorem{lemma}[subsection]{Lemma}

\theoremstyle{definition}
\newtheorem{definition}[subsection]{Definition}
\newtheorem{example}[subsection]{Example}
\newtheorem{exercise}[subsection]{Exercise}
\newtheorem{situation}[subsection]{Situation}

\theoremstyle{remark}
\newtheorem{remark}[subsection]{Remark}
\newtheorem{remarks}[subsection]{Remarks}

\numberwithin{equation}{subsection}

% Macros
%
\def\lim{\mathop{\mathrm{lim}}\nolimits}
\def\colim{\mathop{\mathrm{colim}}\nolimits}
\def\Spec{\mathop{\mathrm{Spec}}}
\def\Hom{\mathop{\mathrm{Hom}}\nolimits}
\def\Ext{\mathop{\mathrm{Ext}}\nolimits}
\def\SheafHom{\mathop{\mathcal{H}\!\mathit{om}}\nolimits}
\def\SheafExt{\mathop{\mathcal{E}\!\mathit{xt}}\nolimits}
\def\Sch{\mathit{Sch}}
\def\Mor{\mathop{\mathrm{Mor}}\nolimits}
\def\Ob{\mathop{\mathrm{Ob}}\nolimits}
\def\Sh{\mathop{\mathit{Sh}}\nolimits}
\def\NL{\mathop{N\!L}\nolimits}
\def\CH{\mathop{\mathrm{CH}}\nolimits}
\def\proetale{{pro\text{-}\acute{e}tale}}
\def\etale{{\acute{e}tale}}
\def\QCoh{\mathit{QCoh}}
\def\Ker{\mathop{\mathrm{Ker}}}
\def\Im{\mathop{\mathrm{Im}}}
\def\Coker{\mathop{\mathrm{Coker}}}
\def\Coim{\mathop{\mathrm{Coim}}}

% Boxtimes
%
\DeclareMathSymbol{\boxtimes}{\mathbin}{AMSa}{"02}

%
% Macros for moduli stacks/spaces
%
\def\QCohstack{\mathcal{QC}\!\mathit{oh}}
\def\Cohstack{\mathcal{C}\!\mathit{oh}}
\def\Spacesstack{\mathcal{S}\!\mathit{paces}}
\def\Quotfunctor{\mathrm{Quot}}
\def\Hilbfunctor{\mathrm{Hilb}}
\def\Curvesstack{\mathcal{C}\!\mathit{urves}}
\def\Polarizedstack{\mathcal{P}\!\mathit{olarized}}
\def\Complexesstack{\mathcal{C}\!\mathit{omplexes}}
% \Pic is the operator that assigns to X its picard group, usage \Pic(X)
% \Picardstack_{X/B} denotes the Picard stack of X over B
% \Picardfunctor_{X/B} denotes the Picard functor of X over B
\def\Pic{\mathop{\mathrm{Pic}}\nolimits}
\def\Picardstack{\mathcal{P}\!\mathit{ic}}
\def\Picardfunctor{\mathrm{Pic}}
\def\Deformationcategory{\mathcal{D}\!\mathit{ef}}
\usepackage{cleveref}
\crefname{originalremark}{original remark}{original remarks}
\Crefname{originalremark}{Original remark}{Original remarks}
\makeatletter
\newlabel{remark-idempotent-symmetry}{{0H9N}{source}{}{sourcepointer-1}{}}
\newlabel{remark-idempotent-symmetry@cref}{{[originalremark][0][]0H9N}{[1][0][]source}{}{}{}}
\makeatother
\crefname{originalremark}{original remark}{original remarks}
\Crefname{originalremark}{Original remark}{Original remarks}
\makeatletter
\newlabel{remark-direct-sum}{{0103}{source}{}{sourcepointer-2}{}}
\newlabel{remark-direct-sum@cref}{{[originalremark][0][]0103}{[1][0][]source}{}{}{}}
\makeatother
\title{An idempotent endomorphism splits its object into the kernel and the kernel of its complement}
\author{The Stacks Project Authors}
\date{stacks-project:homology.tex:local-snapshot-20261012; GNU FDL 1.2 or later}
\begin{document}
\maketitle
\noindent\textit{Editorial note.} All delivered mathematical text is the author's own, copied from the pinned Stacks Project chapter file homology.tex; no statement or proof step was rewritten, reordered or removed.\par\medskip
\noindent\textit{Editorial note.} The proof is the author's complete argument; the statement is of the form 'constructed as in the statement of Lemma ...', so that statement (source lines 583-609) is delivered as a context span with role=context to keep the card self-contained.\par\medskip
\noindent\textit{Editorial note.} Every labelled reference occurring inside the counted spans is carried in assembly.external\_references as a non-load-bearing navigation pointer (same-chapter ones at their definition site, cross-chapter ones at the invocation inside the counted span), so no load-bearing mathematics is replaced by a stub. No bibliography citation occurs inside the counted spans.\par\medskip
\noindent\textit{Editorial note.} The same-name check for this card scanned every staged Stacks chapter of this repository -- all sources/stacks-*/pinned-original-source.tex files, the card's own chapter excluded -- and the labels found are recorded in dedup.same\_name\_cross\_chapter on the card.\par\medskip
\noindent\textit{Editorial note.} The preamble is the chapter's own native preamble excerpt already used for this source: the comment-type reference, slogan and history environments (native preamble lines 8-13), the Xy-pic setup (native preamble lines 14-18) and the native theorem declarations with the mathematical macros (native preamble lines 154-217, the actual last line of the file).\par\medskip
\noindent\textit{Allowed standard prerequisites (editorial).}\begin{itemize}
\item Kernel and cokernel of an idempotent and of its complement. Same chapter, Lemma lemma-idempotent-kernel-cokernel, delivered on this card as a context span. Supplies the four morphisms constructed out of f and g = id - f and the identities p i = id and q j = id that the proof uses.
\item Direct sums described by four morphisms. Same chapter, Remark remark-direct-sum. Supplies the characterization of a direct sum by four morphisms satisfying the identities, which is how the splitting is recognised.
\end{itemize}\medskip
\ifdefined\hypertarget\hypertarget{sourcepointer-1}{}\fi
\noindent\textit{Source reference.} Original remark 0H9N: sources/stacks-homology/pinned-original-source.tex, source lines 575--575. Navigation pointer invoked at proof line 589; the counted span contains the pointer only and the cited result is not reproduced inside it, so the pointer is registered as a non-load-bearing navigation stub in the established form for Stacks material. No bibliography citation is involved.\par\medskip
\ifdefined\hypertarget\hypertarget{sourcepointer-2}{}\fi
\noindent\textit{Source reference.} Original remark 0103: sources/stacks-homology/pinned-original-source.tex, source lines 372--372. Navigation pointer invoked at proof line 663; the counted span contains the pointer only and the cited result is not reproduced inside it, so the pointer is registered as a non-load-bearing navigation stub in the established form for Stacks material. No bibliography citation is involved.\par\medskip
\begin{lemma}
\label{lemma-idempotent-kernel-cokernel}
Let $\mathcal{A}$ be a preadditive category,
$x$ be an object of $\mathcal{A}$,
and $f : x \to x$ be idempotent.
Write $g = \text{id} _ x - f$
and recall by Remark \ref{remark-idempotent-symmetry}
that $f \circ g = g \circ f = 0$.
\begin{enumerate}
\item If $f$ has kernel $i : \Ker (f) \to x$
and $p$ is the unique morphism such that $g = i \circ p$,
then $p \circ i = \text{id} _ {\Ker (f)}$
and $p$ is the cokernel of $f$.
\item If $f$ has cokernel $p : x \to \Coker (f)$
and $i$ is the unique morphism such that $g = i \circ p$,
then $p \circ i = \text{id} _ {\Coker (f)}$
and $i$ is the kernel of $f$.
\item If $g$ has kernel $j : \Ker (g) \to x$
and $q$ is the unique morphism such that $f = j \circ q$,
then $q \circ j = \text{id} _ {\Ker (g)}$
and $q$ is the cokernel of $g$.
\item If $g$ has cokernel $q : x \to \Coker (g)$
and $j$ is the unique morphism such that $f = j \circ q$,
then $q \circ j = \text{id} _ {\Coker (g)}$
and $j$ is the kernel of $g$.
\end{enumerate}
\end{lemma}

\begin{lemma}
\label{lemma-idempotent-splitting}
Let $\mathcal{A}$ be a preadditive category,
$x$ be an object of $\mathcal{A}$,
and $f : x \to x$ be idempotent.
Write $g = \text{id} _ x - f$.
If each of $f$ and $g$ have at least one of a kernel or cokernel
and in turn morphisms $i , j , p , q$ are constructed
as in the statement of Lemma \ref{lemma-idempotent-kernel-cokernel},
then all four (co)kernels exist and
\begin{align*}
x
& \simeq \Ker (f) \oplus \Ker (g) \\
& \simeq \Coker (f) \oplus \Ker (g) \\
& \simeq \Ker (f) \oplus \Coker (g) \\
& \simeq \Coker (f) \oplus \Coker (g)
\end{align*}
with the structure of each direct product given by $i , j , p , q$.
\end{lemma}

\begin{proof}
By Lemma \ref{lemma-idempotent-kernel-cokernel}
$p \circ i = \text{id} _ {\Ker (f)} = \text{id} _ {\Coker (f)}$
and $q \circ j = \text{id} _ {\Ker (g)} = \text{id} _ {\Coker (g)}$.
By construction, $i \circ p + j \circ q = f + g = \text{id} _ x$.
By Remark \ref{remark-direct-sum} the claim follows.
\end{proof}

\end{document}
